\chapter{Introduction}
\label{ch:introduction}

\section{Background}
\label{sec:background}

The building stock represents a major source of global energy demand and a central target of climate policy. Within the energy demand of buildings, residential buildings account for 70\%, while commercial and public buildings account for the remaining 30\% \parencite{iea2025}. In the European Union, buildings account for approximately 40\% of final energy consumption and 36\% of energy-related greenhouse gas emissions \parencite{epbd2024}. These magnitudes make reductions in the energy demand and emissions of existing buildings central to European climate policy. Accordingly, the 2024 recast of the Energy Performance of Buildings Directive establishes the objective of achieving a zero-emission building stock by 2050 \parencite{epbd2024}. Achieving this objective requires reliable information on the existing building stock to support large-scale energy-performance assessment and the prioritisation of renovation measures \parencite{hettinga2023}. The policy challenge therefore translates into an information challenge at the scale of the existing building stock.

\Ac{UBEM} provides a way to assess this performance at the scale of groups of buildings. \ac{UBEM} applies physical models of heat and mass flows in and around buildings to predict operational energy use and indoor and outdoor environmental conditions \parencite{reinhart2016}. As shown in Figure~\ref{fig:ubem_workflow}, a basic \ac{UBEM} workflow combines climate data, building geometry, construction standards, and usage schedules as inputs to simulations of energy use and environmental performance. Although detailed inputs can be collected for individual buildings or small samples, collecting them consistently across larger urban areas is resource intensive and often impractical \parencite{reinhart2016}. The scalability of \ac{UBEM} therefore depends on representing the building stock with information that can be obtained consistently at urban scale.

\begin{figure}[htbp]
	\centering
	\includegraphics[width=\textwidth]{figures/fig/fig_3_1_UBEM_workflow}
	\caption{General workflow for urban building energy modelling. Geometry, building attributes, and archetype information provide inputs for simulation. The results support assessments of energy use, carbon emissions, and energy performance. Based on \textcite{ang2020}.}
	\label{fig:ubem_workflow}
\end{figure}

Building archetypes address this scalability constraint by representing groups of buildings with similar physical and operational characteristics \parencite{reinhart2016,shen_archetype_2024}. They are used extensively in national and regional bottom-up building stock models. Archetype generation consists of two complementary steps: segmentation and characterisation. Segmentation divides the building stock into groups according to characteristics such as building shape, age, use, climate, and technical systems. Characterisation then assigns each archetype a complete set of thermal and operational properties, including construction assemblies, usage patterns, and building systems \parencite{reinhart2016}. Depending on the application scale and selected segmentation parameters, one archetype may represent fewer than 50 or as many as 500,000 buildings \parencite{reinhart2016}. Archetypes therefore reduce building-specific input requirements, but their usefulness depends on assigning each building to an appropriate group.

European building typologies operationalise this grouping process at national, regional, and municipal scales \parencite{loga2016}. Among them, the \ac{TABULA} project provides archetype matrices for 20 European countries that classify residential buildings by construction period and size class \parencite{loga2016}. Figure~\ref{fig:tabula_example} shows the \ac{TABULA} typology structure, including national building stock models and exemplary buildings. Each archetype cell assigns representative U-values to a group of buildings. By linking these representative properties to the distribution of buildings across cells, national typologies provide a basis for estimating the energy consumption of residential building stocks \parencite{loga2016}. Applying \ac{TABULA} at the level of individual buildings therefore depends on reliable building type and construction year information.

\begin{figure}[htbp]
	\centering
	\includegraphics[width=\textwidth]{figures/fig/fig_1_1_tabula_example}
	\caption{TABULA typology concept. Structural information and residential building stock data from 20 European countries support the development of national building stock models and exemplary buildings. Based on the TABULA WebTool \parencite{tabulawebtool}.}
	\label{fig:tabula_example}
\end{figure}

\section{Motivation}
\label{sec:motivation}

However, the attributes required for building-level archetype assignment are not consistently available. Building type and construction year are often incomplete or inaccessible at urban scale \parencite{milojevicdupont2023,dabrock2025}. One potential source is \Ac{OSM}, a collaborative open geographic database that provides building data across European countries. However, its building coverage and attribute availability vary geographically \parencite{milojevicdupont2023}. Local building cadastres, commercial datasets, remote sensing data, and occupant surveys provide alternative sources, but differences in coverage, accessibility, and attribute definitions limit their consistent use across regions \parencite{dabrock2025}. Taken together, these sources do not provide a consistent basis for obtaining building type and construction year information across regions.

Existing research addresses this information gap in two complementary ways. First, studies enrich building records with attributes inferred from census, geospatial, remote sensing, or street view data. These attributes have supported archetype assignment, heat demand modelling, or the enrichment of several building attributes, although the reference sources and validation levels differ across studies \parencite{dabrock2025,wurm2021,liang2025}. Second, other studies map imagery directly to registered energy classes, either alone or together with aerial, geographic, or thermal information \parencite{mayer2023uk,mayer2023dk,liu2025,sun2026,pan_nd}. Building attribute enrichment produces intermediate building attributes that can be inspected and used in subsequent analyses, whereas direct prediction produces an energy class estimate without predicting these attributes or assigning an archetype. Overall, prior attribute enrichment and direct prediction studies establish that both approaches are feasible, but they do not compare building attributes predicted from imagery and energy class predictions with linked reference data for the same buildings \parencite{dabrock2025,wurm2021,liang2025,mayer2023uk,mayer2023dk,liu2025,sun2026,pan_nd}.

The unresolved issue is how building attributes predicted from imagery compare with reference building attributes and how their errors affect \ac{TABULA} assignment and energy class prediction for the same buildings. Addressing this issue requires linked imagery, reference building attributes, reconstructed geometry, and registered energy classes.

The Netherlands provides the linked building data required for this comparison. The \ac{BAG}, the Dutch national register of addresses and buildings, supplies the building identifier and construction year \parencite{kadasterbag}. Through this identifier, each BAG record can be linked to reconstructed geometry from 3DBAG \parencite{threedbagdocs,peters2022,dukai2021} and to residential building type and registered energy class from EP-Online \parencite{rvoeponline}. Mapillary imagery is subsequently associated with the corresponding BAG building \parencite{mapillary}. The resulting dataset links imagery, reference building attributes, reconstructed geometry, and registered energy class for each building. Because EP-Online does not cover every Dutch building \parencite{hettinga2023}, this setting provides both the reference data required for controlled evaluation and a practical case for predicting building attributes where an \ac{EPC} is unavailable.

\section{Research Objectives}
\label{sec:research_objectives}

Building on the preceding gap, this study investigates whether building attributes predicted from \ac{SVI} can support building energy assessment when reference building attributes are unavailable. The evaluation uses the Dutch \ac{TABULA} typology, referred to in this study as TABULA-NL, to examine predicted attributes through archetype assignment and energy class prediction. Building type and construction year determine the TABULA-NL archetype cell, while floor count is evaluated as an additional input to energy class prediction.

Accordingly, this study addresses three research objectives:

\begin{enumerate}
	\item \textbf{Assess the predictive performance of \ac{SVI} for building type, construction year, and floor count.}

	\item \textbf{Quantify how errors in predicted building type and construction year affect TABULA-NL archetype assignment and estimates of the floor-area-normalised transmission heat transfer coefficient.}

	\item \textbf{Determine how the source of building information affects energy class prediction by comparing reference building attributes, building attributes predicted from \ac{SVI}, and imagery alone.}
\end{enumerate}

These objectives evaluate the complete chain from building attributes predicted from \ac{SVI} to archetype parameters and energy class prediction at building level within predefined data partitions.
